\section{Diseñar el controlador lineal}
\subsection{Dificultades encontradas}

Se han experimentado varios inconvenientes durante la simulación. En primer lugar, se encontró una incompatibilidad de variables en el entorno de simulación Simulink, lo cual hizo imposible encontrar una solución adecuada. El entorno de integración numérica utilizado en Simulink no permite la integración por lapsos, lo que limitó aún más las posibilidades de encontrar una solución viable.

Con el objetivo de superar estas dificultades, se decidió trasladar la simulación al entorno de programación MATLAB. Sin embargo, al utilizar la función \textbf{ode45} incorporada en MATLAB, se encontró que esta no permitía trabajar con lapsos de tiempo que brindaran una solución precisa. Esto representó otro obstáculo importante en el proceso de simulación. Para resolver este problema, se tuvo que desarrollar una variación personalizada de la función \textbf{ode45}  en MATLAB. Esta adaptación permitió trabajar con lapsos de tiempo más adecuados para obtener una solución correcta y confiable en la simulación.


\subsection{Resultados PI en configuración estable}

En las primeras pruebas realizadas con un controlador PI en la versión estable del levitador magnético proporcionaron resultados alentadores. Esta elección se basó en la necesidad de analizar el comportamiento del sistema bajo un controlador ampliamente utilizado y estudiado.

La implementación del controlador PI permitió regular la posición del levitador magnético de manera más precisa y estable. Al ajustar los parámetros del controlador proporcional e integral, se pudo lograr un equilibrio entre la respuesta rápida del sistema y la minimización de errores a largo plazo.

La forma del controlador 
\[
C_{(s)} = \frac{0.0108}{s}
\]


\subsection{Resultados PII en configuración estable}

Se diseñó el controlador con doble efecto integral en la configuración inestable del sistema de levitación magnética. 

La forma del controlador 
\[
C_{(s)} = 53\frac{(s+35)(s+39)(s+15)(s+10)}{s^2(s+700)(s+100)}
\]


% \begin{figure}[h]
% \centering
% \begin{tikzpicture}[node distance=2cm, auto]
%     % Define block styles
%     \tikzstyle{block} = [rectangle, draw, fill=blue!20, text width=10em, text centered, rounded corners, minimum height=4em]
%     \tikzstyle{line} = [draw, -latex']
    
%     % Place nodes
%     \node [block] (system) {Condiciones iniciales};
%     \node [block, below of=system] (modeling) {Controlador};
%     \node [block, below of=modeling] (analysis) {Planta};
%     \node [block, below of=analysis] (design) {Control Design};
%     \node [block, below of=design] (implementation) {Implementation};
%     \node [block, below of=implementation] (testing) {Testing};
%     \node [block, below of=testing] (deployment) {Deployment};
    
%     % Draw edges
%     \path [line] (system) -- (modeling);
%     \path [line] (modeling) -- (analysis);
%     \path [line] (analysis) -- (design);
%     \path [line] (design) -- (implementation);
%     \path [line] (implementation) -- (testing);
%     \path [line] (testing) -- (deployment);
% \end{tikzpicture}
% \end{figure}






\begin{figure}[h]
\centering
   \begin{tikzpicture}[node distance=1.5cm, auto]
        % Define block styles
     \tikzstyle{block} = [rectangle, draw, fill=blue!20, text width=6em, text centered, rounded corners, minimum height=4em]
     \tikzstyle{sum} = [circle, draw, fill=blue!20, inner sep=0pt, minimum size=0.6cm]
     \tikzstyle{arrow} = [draw, -latex']


        % Ubicación de los nodos
      \node[block](planta){Planta};
      \node[block, left of=planta, node distance=3cm](controlador){Controlador};
      \node[sum, left of= controlador, node distance=3cm](suma){X};

    % Dibujo de las líneas
   \draw [arrow](suma) -- (controlador); 
   \draw [arrow](controlador) -- (planta); 
   \draw [arrow] (planta.east) -- ++(1,0) |- ($(suma.east) + (1,-1)$) -| (suma.south) node[pos=0.99] {$-$};

   \draw [arrow] ($(suma.west) - (2,0)$) -- (suma.west)node[midway, above] {referencia +};
    \draw [arrow] (planta) -- ++(4,0)node[midway, above] {posición};







    % \draw [arrow] (planta) -- node[anchor=west] {$u$}(suma);

     % \draw [arrow](suma) -- (controlador)
     % \draw [arrow]() -- ()
        
        % \node [sum, left of=controller, node distance=3cm] (sum1) {};
       
        % \node [sum, right of=actuator, node distance=3cm] (sum2) {};
        % \node [block, above of=sum2, node distance=1.5cm] (plant) {Plant};
        
        
        % % Draw arrows
        % \draw [arrow] (analysis) -- (controller);
        % \draw [arrow] (controller) -- node[name=u] {$u$} (actuator);    
        % \draw [arrow] (plant) -| node[pos=0.95] {$-$} (sum1);



       
    \end{tikzpicture}
    \caption{Diagrama de control}
    \label{Esquema de control}
\end{figure}


